\medskip

Parmi les angines, un quart nécessite la prise d'antibiotiques, les autres non.

Afin d'éviter de prescrire inutilement des antibiotiques, les médecins disposent d'un test de diagnostic ayant les caractéristiques suivantes :

\begin{itemize}
\item[$\bullet~~$]lorsque l'angine nécessite la prise d'antibiotiques, le test est positif dans 90\,\% des cas;
\item[$\bullet~~$]lorsque l'angine ne nécessite pas la prise d'antibiotiques, le test est négatif dans 95\,\% des cas.
\end{itemize}





\bigskip

\textbf{Partie 1}

\medskip

Un patient atteint d'angine et ayant subi le test est choisi au hasard. 

On considère les évènements suivants :

\begin{itemize}
\item[$\bullet~~$]$A$ : « le patient est atteint d'une angine nécessitant la prise d'antibiotiques »{} ; 
\item[$\bullet~~$]$T$ : « le test est positif »{} ;
\item[$\bullet~~$]$\overline{A}$ et $\overline{T}$ sont respectivement les évènements contraires de $A$ et $T$.
\end{itemize}



\begin{enumerate}
\item On résume la situation par un arbre pondéré.

\begin{center}
\bigskip
\pstree[treemode=R,nodesepA=0pt,nodesepB=4pt,levelsep=3.5cm,nrot=:U]{\TR{}}
{
\pstree[nodesepA=4pt]{\TR{$A$}\naput{$0,25$}}
{
\TR{$T$}\naput{$0,90$}
\TR{$\overline{T}$}\nbput{\blue $1-0,90=0,10$}
}
\pstree[nodesepA=4pt]{\TR{$\overline{A}$}\nbput{\blue $1-0,25=0,75$}}
{
\TR{$T$}\naput{\blue $1-0,95=0,05$}
\TR{$\overline{T}$}\nbput{$0,95$} 
}
}
\bigskip
\end{center}

$P(A \cap T) = P(A)\times P_{A}(T)=0,25\times 0,90 = 0,225$

\item 
D'après la formule des probabilités totales:

$P(T)= P(A\cap T) + P\left (\overline{A}\cap T\right ) = 0,225 + 0,75\times 0,05 = \np{0,2625}$

\item On choisit un patient ayant un test positif. La probabilité qu'il soit atteint d'une angine nécessitant la prise d'antibiotiques est:

$P_{T}(A)= \dfrac{P(A\cap T)}{P(T)}=\dfrac{0,225}{\np{0,2625}}\approx \np{0,8571}$

\item
	\begin{enumerate}
		\item Les évènements correspondant à un résultat erroné du test sont:  $\overline{A} \cap T$ et $A \cap \overline{T}$
		.
		\item On définit l'évènement $E$ : \og le test fournit un résultat erroné \fg.
		


$P(E) = P\left (\overline{A} \cap T\right )+P\left (A \cap \overline{T}\right )
= 0,25\times 0,10 + 0,75\times 0,05 = \np{0,0625}$

	\end{enumerate}
\end{enumerate}

\bigskip

\textbf{Partie 2}

\medskip

On sélectionne au hasard un échantillon de $n$ patients qui ont été testés.
On admet que l'on peut assimiler ce choix d'échantillon à un tirage avec remise.
On note $X$ la variable aléatoire qui donne le nombre de patients de cet échantillon ayant un test erroné.



\begin{enumerate}
\item On suppose que $n = 50$.
	\begin{enumerate}
		\item 
On a une répétition de 50 épreuves indépendantes et identiques n'ayant que deux issues et dont le succès a pour probabilité $p=\np{0,0625}$; donc la variable aléatoire $X$ qui donne le nombre de succès suit la loi binomiale $\mathcal{B}(n,\: p)$ de paramètres $n = 50$ et $p = \np{0,0625}.$
		
		\item $P(X = 7) = \ds\binom{50}{7}\times \np{0,0625}^{7}\times (1-\np{0,0625})^{50-7}\approx \np{0,0237}$
		
		\item La probabilité qu'il y ait au moins un patient dans l'échantillon dont le test est erroné est:
		
$P(X\geqslant 1) = 1-P(X=0) = 1- \ds\binom{50}{0}\times \np{0,0625}^0\times (1-\np{0,0625})^{50}\approx \np{0,9603}$		
		
	\end{enumerate}	
\item On cherche la taille de l'échantillon faut-il choisir pour que $P(X \geqslant 10)$ soit supérieure à $0,95$.

$P(X \geqslant 10)>0,95 \iff P(X < 10) \leqslant 0,05 \iff P(X \leqslant 9) \leqslant 0,05$

À la calculatrice, par essais successifs, on trouve:

~~\textbullet~~pour $n=247$, $P(X\leqslant 9) \approx \np{0,0514}$;

~~\textbullet~~pour $n=248$, $P(X\leqslant 9) \approx \np{0,0498}$.

Il faut donc un échantillon de taille au moins 248 pour que $P(X\geqslant 10)\geqslant 0,95$.

\end{enumerate}

\bigskip
